\documentclass{article}
\usepackage{amsmath}
\usepackage{amssymb}
\usepackage[utf8]{inputenc}
\usepackage[T1]{fontenc}
\usepackage{geometry}
\usepackage{listings}
\usepackage{color}
\definecolor{codegreen}{rgb}{0,0.6,0}
\definecolor{codegray}{rgb}{0.5,0.5,0.5}
\definecolor{codepurple}{rgb}{0.58,0,0.82}
\definecolor{backcolour}{rgb}{0.95,0.95,0.92}
\lstdefinestyle{mystyle}{
backgroundcolor=\color{backcolour},
commentstyle=\color{codegreen},
keywordstyle=\color{magenta},
numberstyle=\tiny\color{codegray},
stringstyle=\color{codepurple},
basicstyle=\footnotesize,
breakatwhitespace=false,
breaklines=true,
captionpos=b,
keepspaces=true,
numbers=left,
numbersep=5pt,
showspaces=false,
showstringspaces=false,
showtabs=false,
tabsize=2
}
\lstset{style=mystyle}
\geometry{margin=1in}
\title{Analisis Lengkap Kesalahan Implementasi UMP3 pada Framework MSHQC}
\author{Dr. QuantumChem Assistant \ Spesialis Quantum Chemistry Computational Methods}
\date{December 02, 2025}
\begin{document}
\maketitle
\section{Pengantar}
Halo! Saya Dr. QuantumChem Assistant, spesialis dalam metode komputasi kimia kuantum dengan fokus pada framework MSHQC. Dokumen ini menyajikan analisis lengkap mengenai kesalahan dalam implementasi Unrestricted Møller-Plesset Perturbation Theory orde 3 (UMP3) berdasarkan kode yang disediakan. Analisis ini mencakup konteks fisik, rumus matematis, dan identifikasi bug potensial di berbagai bagian kode.
UMP2 Anda sudah benar, dengan koreksi energi negatif dan kecil (-0.000192 Ha untuk Li cc-pVDZ), yang konsisten dengan korelasi dinamis minimal pada atom Li (konfigurasi 1s²2s¹). Namun, UMP3 menunjukkan koreksi E(3) positif dan terlalu besar (+0.133844 Ha), menyebabkan divergensi seri perturbasi. Ini tidak wajar karena E(3) seharusnya negatif dan sekitar 10-20% dari E(2), yakni sekitar -0.00002 hingga -0.00004 Ha.
Analisis ini didasarkan pada referensi utama:
\begin{itemize}
\item Pople et al. (1977), Int. J. Quantum Chem. 11, 149: Formulasi MP3.
\item Bartlett & Silver (1974), Phys. Rev. A 10, 1927: Ekstensi untuk open-shell.
\item Szabo & Ostlund (1996): Transformasi integral dan konvensi notasi.
\end{itemize}
Nilai benchmark dari pengetahuan umum (cross-reference Psi4/PySCF): Untuk Li cc-pVDZ, UHF $\approx$ -7.43242 Ha, UMP2 corr $\approx$ -0.00019 Ha, UMP3 corr $\approx$ -0.00021 Ha (total corr lebih negatif).
\section{Benchmark Nilai Energi}
Untuk validasi, nilai expected untuk Li (spin=1, basis cc-pVDZ):
\begin{align*}
E_{\text{UHF}} &= -7.43242 , \text{Ha} \
E_{\text{UMP2 corr}} &= -0.00019 , \text{Ha} \
E_{\text{UMP3 corr}} &= -0.00021 , \text{Ha} \quad (E(3) \approx -0.00002 , \text{Ha})
\end{align*}
Kode Anda memberikan E(3) positif, menunjukkan sign error atau indexing salah, terutama di spin mixed (αβ).
\section{Analisis Kesalahan di Transformasi Integral}
Bagian ini (di \texttt{eri_transformer.cc} dan \texttt{integrals.cc}) tampak benar, menggunakan notasi physicist secara konsisten. Namun, potensial bug:
\begin{itemize}
\item \textbf{Indexing di \texttt{transform_oovv_mixed}}: Parameter C_occ_A (α occ), C_occ_B (β occ), C_virt_A (α virt), C_virt_B (β virt). Loop:
\begin{lstlisting}[language=C++]
val += C_occ_A(mu, i) * C_occ_B(nu, j) * eri_ao(mu, nu, lam, sig) * C_virt_A(lam, a) * C_virt_B(sig, b);
\end{lstlisting}
Ini menghitung <iα jβ | aα bβ> physicist, benar untuk opposite-spin. Tapi jika \texttt{eri_ao} dari libcint (chemist notation), konfirmasi konvensi:
$$(\mu\nu|\lambda\sigma)_{\text{chemist}} = \langle \mu\lambda | \nu\sigma \rangle_{\text{physicist}}$$
Kode sesuai, tapi cek jika untuk vvvv mixed, indexing (aα bβ | cα dβ) overcount karena loop nested O(N^8) naif, bisa numerik unstable.
\item \textbf{Kesalahan Potensial}: Pada Li (nvir_a=12, nvir_b=13), W_vvvv_ab storage O(nvir^4) $\approx$ 12^2 \times 13^2 $\approx$ 24k elements, tapi jika quarter transform salah (di \texttt{transform_vvvv}), val overflow atau zero.
Rekomendasi: Print sample eri_oovv_ab(0,0,0,0) dan banding dengan PySCF.
\end{itemize}
\section{Analisis Kesalahan di W-Intermediates}
Di \texttt{ump3.cc}, build W menggunakan transform, tapi indexing mungkin salah untuk mixed-spin.
\begin{itemize}
\item \textbf{W_oooo_ab}: <mn|ij> mixed (m n αocc, i j βocc). Kode: \texttt{ERITransformer::transform_oooo_mixed(C_occ_a, C_occ_b)}. Benar, no antisym untuk mixed.
\item \textbf{W_vvvv_ab}: <ab|ef> mixed (a e αvir, b f βvir). Benar, tapi di Li, sum PP ladder dominan jika W besar.
\item \textbf{W_ovov_ab}: <mb|ej> αocc-βvir. Kode: W(m, e, i, b) = eri_oovv_ab(m, i, e, b) = <m i | e b> physicist.
Standar untuk cross PH: - \sum_{m e} \langle m a | e J \rangle t_{i m}^{e B}
\langle m a | e J \rangle = (m e | a J) chemist.
Tapi di kode, term - W(m, e, i, b) * t(m, j, a, e) = - <m i | e b> t m j a e
Ini tidak match < m a | e J > t i m e b (note indices fixed a J vs i b).
\textbf{Bug Utama}: Indexing tensor W_ovov_ab salah. Seharusnya dimensi nocc_a (m), nvir_b (e), nvir_a (a), nocc_b (J), dengan W(m, e, a, J) = eri_oovv_ab(m, J, e, a) = (m J | e a) = <m e | J a>.
Lalu val -= W(m, e, a, j) * t2_ab_1(i, m, e, b)
Ini match the standard term - < m a | e j > t i m e b
Sama untuk W_ovov_ba.
Solusi: Ubah dimensi tensor dan loop build.
\end{itemize}
\section{Analisis Kesalahan di Komputasi T2^{(2)}}
Di \texttt{compute_t2_ab_2nd()}, val dari ladder negative, PH positive. Di Li, jika ladder dominant (HH/PP), val negative, D negative, t2 positive, E(3) = g >0 * t2 >0 = positive (match output).
\textbf{Kesalahan}: PH terms tidak counter ladder karena indexing salah di W_ovov, menyebabkan PH kontribusi kecil.
Rumus benar untuk t_{iJ}^{aB}^{(2)}:
\begin{align*}
t_{iJ}^{aB}^{(2)} = \frac{1}{D_{iJaB}} \Bigg[ &\sum_{k L} \langle k i | L j \rangle t_{k L}^{a B (1)} + \sum_{c D} \langle a B | c D \rangle t_{i J}^{c D (1)} \
&- \sum_{m e} \langle m i | e a \rangle t_{m J}^{e B (1)} - \sum_{m f} \langle m J | f B \rangle t_{i m}^{a f (1)} \
&- \sum_{m e} \langle m a | e J \rangle t_{i m}^{e B (1)} - \sum_{m f} \langle m B | f i \rangle t_{m J}^{a f (1)} \Bigg]
\end{align*}
Kode Anda cocok konseptual, tapi indexing W salah menyebabkan term cross tidak efektif.
\section{Analisis Kesalahan di Komputasi E^{(3)}}
Di \texttt{compute_e3_ab()}, energy += eri_oovv_ab(i, j, a, b) * t2_ab_2(i, j, a, b) (no 0.25, benar untuk mixed).
Jika t2 positive large, E positive.
\section{Kesimpulan dan Rekomendasi}
Kesalahan utama di indexing W_ovov_ab dan ba, menyebabkan PH terms lemah, ladder dominant, val negative large, t2 positive, E(3) positive.
Rekomendasi:

Perbaiki dimensi W_ovov_ab menjadi (nocc_a, nvir_b, nvir_a, nocc_b)
Update build dan loop sesuai.
Test dengan H2 (closed-shell, UMP3 = RMP3).

Untuk debug, upload full \texttt{ump3.cc} jika masih truncated.
\end{document}